\documentclass[12pt]{article}

\usepackage[a4paper,margin=2.35cm]{geometry}
\usepackage{amsmath,amssymb,bm}
\usepackage{booktabs,longtable,array}
\usepackage{pdflscape}
\usepackage{microtype}
\usepackage{xcolor}
\usepackage{hyperref}
\usepackage{url}
\usepackage{fancyhdr}

\hypersetup{
  colorlinks=true,
  linkcolor=blue!50!black,
  citecolor=blue!50!black,
  urlcolor=blue!50!black,
  pdftitle={Integrated demand-side flexibility},
  pdfauthor={CISPO}
}

\setlength{\parindent}{0pt}
\setlength{\parskip}{0.55em}
\renewcommand{\arraystretch}{1.18}
\renewcommand{\thesection}{S6.\arabic{section}}
\renewcommand{\theequation}{S6-\arabic{equation}}
\renewcommand{\thetable}{S6-\arabic{table}}

\pagestyle{fancy}
\setlength{\headheight}{14.5pt}
\fancyhf{}
\fancyhead[L]{\small S6 INTEGRATED DEMAND-SIDE FLEXIBILITY}
\fancyhead[R]{\small\thepage}
\renewcommand{\headrulewidth}{0pt}

\title{\textbf{S6 Integrated demand-side flexibility}}
\author{}
\date{}

\begin{document}

\section*{S6 Integrated demand-side flexibility}
\addcontentsline{toc}{section}{S6 Integrated demand-side flexibility}

We represent demand-side flexibility as a set of energy-service constraints
rather than as an unconstrained reduction in electricity demand. The
reference system retains the reconstructed hourly demand profiles, whereas
the flexibility case allows selected heating, cooling and electric-vehicle
(EV) services to be shifted in time and permits a limited amount of
vehicle-to-grid (V2G) discharge. Both cases have identical annual electricity
demand, generation and network technologies, security constraints and
economic boundaries. Flexible services incur explicit enabling and
activation costs, and their contribution to resource adequacy is restricted
to the derated power that remains deliverable throughout a four-hour peak
window. This formulation allows demand flexibility to reduce operating costs
and, where physically justified, firm generation requirements without
treating flexible demand as a free or perfectly reliable resource.

\section{Demand reconstruction and flexible-service envelopes}

Hourly electricity demand is reconstructed for each province and decomposed
into residual, space-heating, space-cooling and EV-charging components:
\begin{equation}
  L^{0}_{p,t}
  =
  L^{\mathrm{res}}_{p,t}
  +L^{\mathrm{heat},0}_{p,t}
  +L^{\mathrm{cool},0}_{p,t}
  +L^{\mathrm{EV},0}_{p,t},
  \label{eq:loaddecomposition}
\end{equation}
where $p$ and $t$ denote province and hour, respectively. The decomposition is
closed by construction: heating, cooling and EV charging are internal
components of total demand rather than additions to it. Historical hourly
profiles are anchored to provincial electricity demand, with the Spring
Festival correction modifying intrayear shape while preserving annual
electricity consumption. Future residual demand is obtained by subtracting
projected heating, cooling and EV electricity from provincial annual demand
and distributing the remainder over the retained 8,760-hour residual-load
profile.

Heating and cooling loads are derived from population- and electricity-
weighted ERA5 meteorology \cite{Hersbach2020ERA5}. Temperature, solar
irradiance, wind speed and humidity are combined in a building-adjusted
temperature index with a 48-hour exponentially weighted memory, after which
region-specific heating- and cooling-degree thresholds are applied. For each
province and hour, the response of the reconstructed thermal load to a
$\pm1\,^{\circ}\mathrm{C}$ comfort perturbation defines the maximum upward
and downward modulation. These weather- and load-dependent envelopes prevent
the optimization from creating thermal storage independently of the
underlying end use. The use of an equivalent energy state to represent
aggregated building flexibility is consistent with established
characterizations of thermal mass and controlled heating, ventilation and
air-conditioning loads
\cite{IEAEBC67,Stinner2016,Klein2017,HuangWu2019,Chen2023Cooling}.

The exogenous EV-charging profile is constructed from provincial new-energy
vehicle stocks, monthly electricity use per vehicle and normalized intraday
charging probabilities. For planning year $y$,
\begin{equation}
  L^{\mathrm{EV},0}_{p,h,y}
  =
  \frac{
    N^{\mathrm{EV}}_{p,y}\,
    e_{m(h)}\,
    \pi_h
  }{1000},
  \label{eq:evbaseline}
\end{equation}
where $N^{\mathrm{EV}}_{p,y}$ is the provincial vehicle stock,
$e_{m(h)}$ is daily charging electricity per vehicle in month $m$, and
$\pi_h$ is the normalized hourly charging probability. Coordinated charging
does not alter vehicle stocks, annual mobility energy or the uncontrolled
charging reference. Instead, a prescribed share of this reference is made
schedulable within power, energy and departure-service limits.

The central case enrolls 25\% of heating demand, 20\% of cooling demand and
15\% of EV charging in flexible-service contracts. These fractions scale the
physical envelopes rather than the underlying demand. They are treated as
scenario parameters because national demand-response measures and local
air-conditioning programmes establish the institutional feasibility of
contracted flexibility, but do not provide nationally representative
participation rates
\cite{NDRCDemand2023,ShanghaiHVAC2025,AnhuiDR2025,ZhejiangHVAC2025}.
Observed responses in Chinese coordinated-charging programmes similarly
support a non-zero schedulable EV share while revealing substantial
heterogeneity in participation and price response
\cite{Liu2021OrderlyCharging,Zhao2026Charging}.

\section{Joint thermal-load and electric-vehicle flexibility model}

For thermal service $s\in\{\mathrm{heat},\mathrm{cool}\}$, the optimization
chooses contracted capacity $K^s_p$, upward modulation $P^{s,+}_{p,t}$ and
downward modulation $P^{s,-}_{p,t}$. The hourly decisions are bounded by the
comfort-derived envelopes, $\overline P^{s,+}_{p,t}$ and
$\overline P^{s,-}_{p,t}$, and by the contracted power available at that
hour, $a^s_{p,t}K^s_p$:
\begin{align}
  0\le P^{s,+}_{p,t}
  &\le
  \min\!\left\{\overline P^{s,+}_{p,t},a^s_{p,t}K^s_p\right\},\\
  0\le P^{s,-}_{p,t}
  &\le
  \min\!\left\{\overline P^{s,-}_{p,t},a^s_{p,t}K^s_p\right\}.
  \label{eq:thermalpower}
\end{align}
An equivalent thermal-service state $S^s_{p,t}$ tracks the intertemporal
consequence of pre-heating or pre-cooling:
\begin{equation}
  S^s_{p,t}
  =
  \rho^s_p S^s_{p,t-1}
  +\eta^{s,+}_p P^{s,+}_{p,t}
  -\frac{P^{s,-}_{p,t}}{\eta^{s,-}_p},
  \qquad
  0\le S^s_{p,t}\le d^s_p K^s_p.
  \label{eq:thermalstate}
\end{equation}
Here, $\rho^s_p$ is hourly service retention and $d^s_p$ is the maximum
energy-to-power duration. The state is non-negative and periodic over the
modelled year, precluding terminal comfort debt. The resulting electrical
load is
$\widetilde L^s_{p,t}=L^{s,0}_{p,t}+P^{s,+}_{p,t}-P^{s,-}_{p,t}$.
Central hourly retention is 0.94 for heating and 0.92 for cooling, with
corresponding service durations of 4 and 5 hours.

EV flexibility is represented by a single aggregate mobility-service state
rather than by superimposing a second battery on coordinated charging. A
fraction $\alpha^{\mathrm{V1G}}$ of the uncontrolled EV load is schedulable,
whereas the remainder is fixed:
\begin{equation}
  L^{\mathrm{EV,flex},0}_{p,t}
  =
  \alpha^{\mathrm{V1G}}L^{\mathrm{EV},0}_{p,t},
  \qquad
  L^{\mathrm{EV,fixed}}_{p,t}
  =
  (1-\alpha^{\mathrm{V1G}})L^{\mathrm{EV},0}_{p,t}.
\end{equation}
Let $P^{\mathrm{ch}}_{p,t}$ and $P^{\mathrm{out}}_{p,t}$ denote flexible
charging and V2G discharge, respectively. The fleet-service state evolves as
\begin{equation}
  E^{\mathrm{EV}}_{p,t}
  =
  E^{\mathrm{EV}}_{p,t-1}
  +\eta^{\mathrm{ch}}P^{\mathrm{ch}}_{p,t}
  -D^{\mathrm{mob}}_{p,t}
  -\frac{P^{\mathrm{out}}_{p,t}}{\eta^{\mathrm{dis}}},
  \label{eq:evstate}
\end{equation}
where $D^{\mathrm{mob}}_{p,t}$ is exogenous driving-energy withdrawal. The
state is bounded by the hourly fleet-energy envelope and by a minimum
departure-energy requirement, and it is periodic over the modelled year.
Charging and discharge are further constrained by exogenous hourly power
envelopes and endogenous contracted capacities:
\begin{align}
  0\le P^{\mathrm{ch}}_{p,t}
  &\le
  \min\!\left\{\overline P^{\mathrm{ch}}_{p,t},
  a^{\mathrm{EV}}_{p,t}K^{\mathrm{V1G}}_p\right\},\\
  0\le P^{\mathrm{out}}_{p,t}
  &\le
  \min\!\left\{\overline P^{\mathrm{out}}_{p,t},
  a^{\mathrm{EV}}_{p,t}K^{\mathrm{V2G}}_p\right\},\\
  0\le K^{\mathrm{V2G}}_p
  &\le K^{\mathrm{V1G}}_p.
  \label{eq:evpower}
\end{align}
Thus, bidirectional participation is nested within the coordinated-charging
contract. It cannot exceed either the participating fleet or its hourly
power and energy availability. Charge and discharge efficiencies are both
0.94 in the central case.

The national sum of contracted V2G power is capped at 10, 20, 30 and 40 GW in
2030, 2040, 2050 and 2060, respectively. The 2030 value is anchored to
China's policy objective of establishing bidirectional interaction at the
ten-gigawatt scale, and the subsequent values are explicit planning
extrapolations rather than forecasts. The national vehicle--grid policy and
the first large-scale pilot programme, covering nine cities and 30 projects,
provide the deployment context
\cite{NDRCVGI2023,NDRCVGIPilot2025}. The aggregate representation deliberately
retains conservative service constraints because city-scale evidence shows
that mobility, dwell time, connection, compensation, range requirements and
battery degradation jointly determine usable V2G potential
\cite{IEAV2G2026,Li2025V2G}.

\section{System integration, economic accounting and capacity adequacy}

The hourly load entering the power balance is
\begin{align}
  L^{\mathrm{eff}}_{p,t}
  &=
  L^{\mathrm{res}}_{p,t}
  +\widetilde L^{\mathrm{heat}}_{p,t}
  +\widetilde L^{\mathrm{cool}}_{p,t} \nonumber\\
  &\quad
  +L^{\mathrm{EV,fixed}}_{p,t}
  +P^{\mathrm{ch}}_{p,t}
  -P^{\mathrm{out}}_{p,t}.
  \label{eq:effectiveload}
\end{align}
This effective load is used consistently in hourly power balance and in
load-dependent operating-security requirements. Flexible demand receives no
additional reserve credit, which avoids using the same deliverable power
simultaneously for resource adequacy and operating reserves.

All flexible services incur real-resource costs. Thermal flexibility pays an
annual enabling cost and an activation cost on both upward and downward
modulation. Coordinated EV charging pays an annual control and aggregation
cost and an inconvenience cost on charging energy displaced from its
uncontrolled hour. To prevent V2G replenishment from being misclassified as
coordinated-charging activity, relocated energy is defined by
\begin{equation}
  R^{\mathrm{V1G}}_{p,t}
  \ge
  L^{\mathrm{EV,flex},0}_{p,t}-P^{\mathrm{ch}}_{p,t},
  \qquad
  R^{\mathrm{V1G}}_{p,t}\ge0.
  \label{eq:relocation}
\end{equation}
V2G additionally pays for availability, annualized bidirectional charging
infrastructure, owner participation per unit of discharge and battery
degradation per unit of discharge. The incremental flexibility cost is
\begin{align}
 C^{\mathrm{flex}}
 &=
 \sum_{p,s}c^{\mathrm{ena}}_s K^s_p
 +10^{-3}\sum_{p,t,s\in\{\mathrm{heat},\mathrm{cool}\}}
 c^{\mathrm{act}}_s
 \left(P^{s,+}_{p,t}+P^{s,-}_{p,t}\right) \nonumber\\
 &\quad
 +10^{-3}\sum_{p,t}c^{\mathrm{reloc}}_{\mathrm{V1G}}
 R^{\mathrm{V1G}}_{p,t}
 +\sum_p c^{\mathrm{infra}}_{\mathrm{V2G}}
 K^{\mathrm{V2G}}_p \nonumber\\
 &\quad
 +10^{-3}\sum_{p,t}
 \left(c^{\mathrm{owner}}_{\mathrm{V2G}}
 c^{\mathrm{deg}}_{\mathrm{V2G}}\right)
 P^{\mathrm{out}}_{p,t}.
  \label{eq:flexcost}
\end{align}
Costs are expressed in 2025 constant CNY. Time-of-use tariffs and public
subsidies are treated as transfers and are not added to social system cost;
metering, control, aggregation, infrastructure, inconvenience and battery
wear are treated as real resource costs. Charging electricity and conversion
losses remain endogenous in the system energy balance and are therefore not
charged a second time. This accounting follows the broader evidence that the
system value of demand flexibility depends on non-zero enabling,
interoperability and participation costs
\cite{IEADemandFlex2025,Wan2026V2G,Sagaria2025V2G}.

Resource adequacy is assessed against the immutable annual provincial peak,
rather than against an endogenous peak that could be reduced without a
deliverability test:
\begin{equation}
  C^{\mathrm{credit}}_p
  +\sum_s F^{s,\mathrm{firm}}_p
  \ge
  (1+m)\max_t L^{0}_{p,t},
  \label{eq:adequacy}
\end{equation}
where $C^{\mathrm{credit}}_p$ is supply-side accredited capacity, $m$ is the
planning margin and $F^{s,\mathrm{firm}}_p$ is the accredited contribution of
flexible service $s$. For each province, the deliverability window comprises
the annual peak hour, the preceding hour and the following two hours. Thermal
and coordinated-charging credits are bounded by contracted power and the
minimum reduction envelope throughout this four-hour window. V2G credit is
additionally bounded by the minimum fleet energy in the window divided by
four hours:
\begin{align}
 F^{s,\mathrm{firm}}_p
 &\le
 \delta_s\min\left\{
 K^s_p,\min_{t\in\mathcal W_p}\overline P^{s,-}_{p,t}
 \right\},
 &&s\in\{\mathrm{heat},\mathrm{cool}\},\\
 F^{\mathrm{V1G},\mathrm{firm}}_p
 &\le
 \delta_{\mathrm{V1G}}\min\left\{
 K^{\mathrm{V1G}}_p,
 \min_{t\in\mathcal W_p}L^{\mathrm{EV,flex},0}_{p,t}
 \right\},\\
 F^{\mathrm{V2G},\mathrm{firm}}_p
 &\le
 \delta_{\mathrm{V2G}}\min\left\{
 K^{\mathrm{V2G}}_p,
 \min_{t\in\mathcal W_p}\overline P^{\mathrm{out}}_{p,t},
 \frac{\min_{t\in\mathcal W_p}\overline E^{\mathrm{fleet}}_{p,t}}{4}
 \right\}.
  \label{eq:firmcredit}
\end{align}
Central derating factors are 0.60 for heating and cooling and 0.50 for
coordinated charging and V2G. This approach preserves the capacity value of
demand flexibility while preventing a one-hour reduction in effective load
from being interpreted as perfectly firm capacity. It is consistent with the
principle that capacity accreditation should reflect sustained,
system-specific deliverability
\cite{NREL2024RA,ESIG2023Capacity,Mantegna2024RA}. Demand response and storage
capacity value have also been incorporated in Chinese expansion-planning
research \cite{Huang2023DRCapacity}; however, the adopted derating factors are
engineering assumptions, not empirically estimated provincial ELCC values.
Regulatory ELCC treatments are used only as methodological comparators
\cite{CPUC2026IRP}.

\clearpage
\begin{landscape}
\section{Parameter evidence, sensitivities and interpretation}

Table \ref{tab:flexparameters} summarizes the central assumptions and the
prespecified ranges used to evaluate structural uncertainty. The parameter
set combines reproducible end-use envelopes with Chinese policy and
behavioural evidence, peer-reviewed technology studies and international
resource-adequacy guidance. Policy targets and compensation rules are used to
establish feasibility and scale, rather than being transferred directly into
national social-cost parameters.

\small
\begin{longtable}{
  >{\raggedright\arraybackslash}p{5.1cm}
  >{\centering\arraybackslash}p{2.7cm}
  >{\centering\arraybackslash}p{3.2cm}
  >{\raggedright\arraybackslash}p{9.2cm}}
\caption{Central assumptions and prespecified sensitivity ranges for
integrated demand-side flexibility. Monetary values are in 2025 constant
CNY.}
\label{tab:flexparameters}\\
\toprule
Parameter & Central value & Sensitivity range & Evidence and interpretation\\
\midrule
\endfirsthead
\multicolumn{4}{l}{\small Table \thetable{} continued}\\
\toprule
Parameter & Central value & Sensitivity range & Evidence and interpretation\\
\midrule
\endhead
\midrule
\multicolumn{4}{r}{Continued on next page}\\
\endfoot
\bottomrule
\endlastfoot
Heating and cooling enrollment & 0.25; 0.20 & 0.10--0.40 &
Applied to weather- and comfort-derived envelopes. Chinese demand-response
rules and local air-conditioning programmes establish feasibility but not a
national participation rate
\cite{NDRCDemand2023,ShanghaiHVAC2025,AnhuiDR2025,ZhejiangHVAC2025}.\\
Heating and cooling hourly retention & 0.94; 0.92 &
0.90--0.98; 0.85--0.97 &
Equivalent-service representation informed by building thermal-flexibility
studies \cite{Stinner2016,HuangWu2019,Chen2023Cooling}.\\
Heating and cooling duration & 4 h; 5 h & 2--8 h; 2--10 h &
Energy-to-power limit for thermal service, evaluated across the range reported
for controlled HVAC and thermal mass \cite{IEAEBC67,Klein2017}.\\
Thermal enabling cost & 40 CNY/kW-y & 10--120 CNY/kW-y &
Control, metering, aggregation and availability; treated as a real-resource
cost \cite{IEADemandFlex2025,NDRCDemand2023}.\\
Thermal activation cost & 300 CNY/MWh & 100--1000 CNY/MWh &
Service activation and comfort cost. Local compensation schedules provide
scale comparators but are not added as fiscal transfers
\cite{AnhuiDR2025,ZhejiangHVAC2025}.\\
Coordinated-charging enrollment & 0.15 & 0.10--0.25 &
Share of uncontrolled EV charging that can be rescheduled; informed by
Chinese pilot participation and price-response evidence
\cite{Liu2021OrderlyCharging,Zhao2026Charging}.\\
V2G participation envelope & 0.10 & 0.05--0.30 &
Scales the provincial discharge envelope; does not represent an hourly
plug-in probability \cite{NDRCVGI2023,Li2025V2G}.\\
Charging and discharge efficiencies & 0.94; 0.94 & 0.90--0.97 &
Grid-to-fleet and fleet-to-grid conversion efficiencies
\cite{IEAV2G2026,Li2025V2G}.\\
Coordinated-charging power ratio and service duration & 2.0; 24 h &
1.5--3.0; 12--36 h &
Aggregate power and energy limits evaluated as scenario uncertainty; the
service state is not interpreted as nameplate vehicle-battery capacity.\\
Coordinated-charging enabling and relocation costs &
60 CNY/kW-y; 300 CNY/MWh &
20--180 CNY/kW-y; 100--400 CNY/MWh &
Control and aggregation cost plus participant inconvenience on charging
energy displaced once \cite{Liu2021OrderlyCharging,Zhao2026Charging}.\\
V2G availability and infrastructure costs &
60; 40 CNY/kW-y &
20--180; 20--100 CNY/kW-y &
Incremental availability and annualized bidirectional equipment/platform
costs \cite{IEAV2G2026,Wan2026V2G}.\\
V2G owner and degradation costs &
150; 400 CNY/MWh &
100--300; 250--500 CNY/MWh &
Participation compensation and battery wear are represented separately
\cite{Wan2026V2G,Sagaria2025V2G}.\\
Thermal firm-capacity derating & 0.60 & 0.30--0.80 &
Applied only to four-hour deliverable reduction; not a China-specific ELCC
estimate \cite{NREL2024RA,ESIG2023Capacity,Mantegna2024RA}.\\
Coordinated-charging and V2G firm-capacity derating & 0.50; 0.50 &
0.20--0.70 &
Applied to four-hour power and energy deliverability; not a China-specific
ELCC estimate \cite{Huang2023DRCapacity,CPUC2026IRP}.\\
National V2G contracted-power ceiling &
10/20/30/40 GW & 5--10/10--30/15--50/20--70 GW &
2030 is anchored to China's ten-gigawatt-scale objective; later values are
planning extrapolations \cite{NDRCVGI2023,NDRCVGIPilot2025}.\\
\end{longtable}
\end{landscape}

The integrated formulation is intended to identify mechanisms rather than to
assign a deterministic forecast to demand-side participation. Its results
should therefore be interpreted jointly with the low and high assumptions in
Table \ref{tab:flexparameters}. In particular, a reduction in generation
capacity is attributable to demand-side flexibility only when the associated
four-hour firm-credit constraint is binding and the full enabling,
activation, infrastructure and degradation costs are included. A reduction
in the unconstrained effective peak alone is not sufficient evidence of firm
capacity substitution.

Several limitations remain. Building flexibility is represented by
province-level comfort envelopes rather than building-by-building thermal
models. EV mobility and connection are aggregated, and the model does not
resolve individual trip chains, charger locations or distribution-network
constraints. The four-hour accreditation window is a transparent
deterministic approximation and does not replace a probabilistic
loss-of-load or ELCC assessment. Finally, periodic service states are
appropriate for an annual planning model but should not be interpreted as
evidence about behaviour beyond the represented horizon. These limitations
motivate sensitivity analysis of participation, duration, costs and capacity
derating, and constrain scientific interpretation to complete annual
simulations.

{\small\sloppy\begin{thebibliography}{99}

\bibitem{Hersbach2020ERA5}
Hersbach H, Bell B, Berrisford P, et al.
\newblock The ERA5 global reanalysis.
\newblock \emph{Quarterly Journal of the Royal Meteorological Society},
146:1999--2049, 2020.
\newblock \url{https://doi.org/10.1002/qj.3803}.

\bibitem{IEAEBC67}
IEA Energy in Buildings and Communities Programme.
\newblock \emph{Annex 67: Energy Flexible Buildings}.
\newblock International Energy Agency, 2019.
\newblock \url{https://www.iea-ebc.org/projects/project?AnnexID=67}.

\bibitem{Stinner2016}
Stinner S, Huchtemann K, Müller D.
\newblock Quantifying the operational flexibility of building energy systems
with thermal energy storages.
\newblock \emph{Applied Energy}, 181:140--154, 2016.
\newblock \url{https://doi.org/10.1016/j.apenergy.2016.08.055}.

\bibitem{Klein2017}
Klein K, Herkel S, Henning H-M, Felsmann C.
\newblock Load shifting using the heating and cooling system of an office
building: quantitative potential evaluation for different flexibility and
storage options.
\newblock \emph{Applied Energy}, 203:917--937, 2017.
\newblock \url{https://doi.org/10.1016/j.apenergy.2017.06.073}.

\bibitem{HuangWu2019}
Huang S, Wu D.
\newblock Validation on aggregate flexibility from residential air
conditioning systems for building-to-grid integration.
\newblock \emph{Energy and Buildings}, 200:58--67, 2019.
\newblock \url{https://doi.org/10.1016/j.enbuild.2019.07.043}.

\bibitem{Chen2023Cooling}
Chen Q, Kuang Z, Liu X, Zhang T.
\newblock Techno-economic comparison of cooling storage and battery for
electricity flexibility.
\newblock \emph{Energy Conversion and Management}, 289:117183, 2023.
\newblock \url{https://doi.org/10.1016/j.enconman.2023.117183}.

\bibitem{NDRCDemand2023}
National Development and Reform Commission of China.
\newblock \emph{Power Demand-Side Management Measures (2023 Edition)}, 2023.
\newblock \url{https://www.ndrc.gov.cn/xxgk/zcfb/ghxwj/202309/t20230927_1360902.html}.

\bibitem{ShanghaiHVAC2025}
Shanghai Municipal People's Government.
\newblock \emph{Notice on Building Air-Conditioning Load Regulation Capability
in 2025}, 2025.
\newblock \url{https://www.shanghai.gov.cn/qtsj/20250617/9db446ddad1744b6af10b6f7205d8aab.html}.

\bibitem{AnhuiDR2025}
Anhui Provincial Energy Administration.
\newblock \emph{Anhui Provincial Electricity Demand Response Implementation
Rules}, 2025.
\newblock \url{https://www.ahjd.gov.cn/Jczwgk/show/3553986.html}.

\bibitem{ZhejiangHVAC2025}
Quzhou Municipal People's Government.
\newblock \emph{Da Kechuang Special Policy Implementation Rules
(2025 Edition)}, 2025.
\newblock \url{https://zjjcmspublic.oss-cn-hangzhou-zwynet-d01-a.internet.cloud.zj.gov.cn/jcms_files/jcms1/web3084/site/attach/0/adde479fede6497d921889690afad168.pdf}.

\bibitem{IEADemandFlex2025}
International Energy Agency.
\newblock \emph{The Value of Demand Flexibility}.
\newblock IEA, Paris, 2025.
\newblock \url{https://www.iea.org/reports/the-value-of-demand-flexibility}.

\bibitem{NDRCVGI2023}
National Development and Reform Commission of China and National Energy
Administration.
\newblock \emph{Implementation Opinions on Strengthening Vehicle--Grid
Integration}, 2023.
\newblock \url{https://zfxxgk.nea.gov.cn/2023-12/13/c_1310758710.htm}.

\bibitem{NDRCVGIPilot2025}
National Development and Reform Commission of China and National Energy
Administration.
\newblock \emph{First Batch of Large-Scale Vehicle--Grid Interaction Pilots},
2025.
\newblock \url{https://www.ndrc.gov.cn/xxgk/zcfb/tz/202504/t20250402_1396958.html}.

\bibitem{Liu2021OrderlyCharging}
Liu X, Li X, Chen H, Xu H.
\newblock Aggregate electric vehicles in demand side for ancillary service.
\newblock In: \emph{34th International Electric Vehicle Symposium and
Exhibition (EVS34)}, Nanjing, 2021.
\newblock \url{https://www.citelec.org/EVS34/download.php?f=papers/evs34-5010511.pdf}.

\bibitem{Zhao2026Charging}
Zhao X, Chen H, Qiu Y, Deng Y.
\newblock Impact of charging price subsidies on the charging behavior of
heterogeneous electric vehicle drivers: evidence from Beijing.
\newblock \emph{Journal of the Association of Environmental and Resource
Economists}, 13(3):713--754, 2026.
\newblock \url{https://doi.org/10.1086/739399}.

\bibitem{IEAV2G2026}
International Energy Agency.
\newblock \emph{Vehicle-to-Grid Technology}.
\newblock IEA, Paris, 2026.
\newblock \url{https://www.iea.org/reports/vehicle-to-grid-technology}.

\bibitem{Li2025V2G}
Li K, Li X, Xiong Z, Tao S, Zhao G, Jiang Y, Qi H, Zhang Y.
\newblock Unlocking vehicle-to-grid potential of load shifting in China's
megacities considering comprehensive real-world behaviors.
\newblock \emph{Nature Communications}, 16:10087, 2025.
\newblock \url{https://doi.org/10.1038/s41467-025-65073-8}.

\bibitem{Wan2026V2G}
Wan M, Xia L, Huo Y, Gong Y, Geng G, Jiang Q.
\newblock Case study of vehicle-to-grid incentive design under multiple
scenarios: a policy-making perspective.
\newblock \emph{IET Generation, Transmission \& Distribution},
20(1):e70254, 2026.
\newblock \url{https://doi.org/10.1049/gtd2.70254}.

\bibitem{Sagaria2025V2G}
Sagaria S, van der Kam M, Boström C.
\newblock Vehicle-to-grid impact on battery degradation and estimation of V2G
economic compensation.
\newblock \emph{Applied Energy}, 377:124546, 2025.
\newblock \url{https://doi.org/10.1016/j.apenergy.2024.124546}.

\bibitem{NREL2024RA}
National Renewable Energy Laboratory.
\newblock \emph{Explained: Fundamentals of Power Grid Reliability and Clean
Electricity}. NREL/TP-6A40-85880, 2024.
\newblock \url{https://docs.nrel.gov/docs/fy24osti/85880.pdf}.

\bibitem{ESIG2023Capacity}
Energy Systems Integration Group.
\newblock \emph{Ensuring Efficient Reliability: New Design Principles for
Capacity Accreditation}, 2023.
\newblock \url{https://www.esig.energy/wp-content/uploads/2023/02/ESIG-Design-principles-capacity-accreditation-report-2023.pdf}.

\bibitem{Mantegna2024RA}
Mantegna G, Huang Z, Van Caelenberg G, Frew B, Lynch M, O'Malley M.
\newblock Maintaining reliability while navigating unprecedented uncertainty:
a synthesis of and guide to advances in electric sector resource adequacy.
\newblock \emph{arXiv:2412.00533}, 2024.
\newblock \url{https://doi.org/10.48550/arXiv.2412.00533}.

\bibitem{Huang2023DRCapacity}
Huang Y, Zhang Y, Xia Z, Wang H, Wu M, Wang N, Chen Q, Zhu T, Chen X.
\newblock Power system planning considering demand response resources and
capacity value of energy storage.
\newblock \emph{Journal of Shanghai Jiao Tong University},
57(4):432--441, 2023.
\newblock \url{https://doi.org/10.16183/j.cnki.jsjtu.2021.477}.

\bibitem{CPUC2026IRP}
California Public Utilities Commission.
\newblock \emph{Inputs and Assumptions for 2026 Integrated Resource Planning},
2026.
\newblock \url{https://www.cpuc.ca.gov/industries-and-topics/electrical-energy/electric-power-procurement/long-term-procurement-planning/2024-26-irp-cycle-events-and-materials}.

\end{thebibliography}
}

\end{document}
